\documentclass[border=10pt]{standalone}
\usepackage{tikz,ctex}
\usetikzlibrary{positioning}

\begin{document}

\begin{tikzpicture}[
    % 全局设置
    node distance = 1.6cm and 1.4cm,
    every node/.style = {
        draw,
        rounded corners,
        thick,
        align = center,
        font = \sffamily\small,
        minimum width = 2.0cm,
        minimum height = 0.9cm
    },
    % 样式定义
    centerStyle/.style = {
        fill = blue!45,
        text = white,
        font = \sffamily\bfseries\large,
        line width = 2pt,
        minimum size = 2.5cm,
        circle
    },
    branchTitle/.style = {
        font = \sffamily\bfseries,
        line width = 1.5pt
    },
    subNode/.style = {
        fill = white,
        font = \sffamily\small,
        minimum width = 2.0cm,
        minimum height = 0.8cm
    },
    leafNode/.style = {
        fill = white,
        font = \sffamily\footnotesize,
        minimum width = 1.8cm,
        minimum height = 0.7cm
    },
    arrowStyle/.style = {
        ->,
        thick,
        >=stealth,
        shorten >=2pt
    }
]

    % ================= 1. 中心节点 =================
    \node[centerStyle] (center) {7.2.梯度下降优化 \\ Gradient Descent};

    % ================= 2. 上方分支：梯度信息价值 (红) =================
    \node[branchTitle, fill=red!15, above=of center] (value) {梯度信息价值 \\ Value of Gradient};

    \node[subNode, above=of value, xshift=-6cm] (iterative) {迭代数值方法 \\ Iterative Methods};
    \node[subNode, above=of value, xshift=-2cm] (info) {信息效率 \\ Information Efficiency};
    \node[subNode, above=of value, xshift=2cm] (complexity) {复杂度对比 \\ Complexity Comparison};
    \node[subNode, above=of value, xshift=6cm] (backprop) {反向传播 \\ Backpropagation};

    % ================= 3. 右侧分支：批量梯度下降 (蓝) =================
    \node[branchTitle, fill=orange!15, right=of center] (batch) {批量梯度下降 \\ Batch Gradient Descent};

    \node[subNode, right=of batch, yshift=3cm] (fullbatch) {全数据集 \\ Full Data Set};
    \node[subNode, right=of batch, yshift=1cm] (lr) {学习率 $\eta$ \\ Learning Rate};
    \node[subNode, right=of batch, yshift=-1cm] (steepest) {最速下降 \\ Steepest Descent};
    \node[subNode, right=of batch, yshift=-3cm] (inefficient) {大数据集低效 \\ Inefficient for Big Data};

    % ================= 4. 下方分支：随机梯度下降 (绿) =================
    \node[branchTitle, fill=green!15, below=of center] (sgd) {随机梯度下降 \\ SGD};

    \node[subNode, below=of sgd, xshift=-6cm] (single) {单样本更新 \\ Single Sample};
    \node[subNode, below=of sgd, xshift=-3cm] (epoch) {训练轮次 \\ Training Epoch};
    \node[subNode, below=of sgd, xshift=0cm] (online) {在线梯度下降 \\ Online GD};
    \node[subNode, below=of sgd, xshift=3cm] (redundancy) {冗余数据高效 \\ Handles Redundancy};
    \node[subNode, below=of sgd, xshift=7cm] (escape) {逃离局部极小 \\ Escape Local Minima};

    % ================= 5. 左侧分支：小批量与初始化 (紫) =================
    \node[branchTitle, fill=purple!15, left=of center] (mini) {小批量与初始化 \\ Mini-batch \& Init};

    \node[subNode, left=of mini, yshift=4cm] (minibatch) {小批量 \\ Mini-batch};
    \node[subNode, left=of mini, yshift=2cm] (noise) {噪声梯度 \\ Noisy Gradient};
    \node[subNode, left=of mini, yshift=0cm] (shuffle) {数据打乱 \\ Data Shuffling};
    \node[subNode, left=of mini, yshift=-2cm] (symmetry) {对称性破缺 \\ Symmetry Breaking};
    \node[subNode, left=of mini, yshift=-4cm] (he) {He 初始化 \\ He Initialization};

    % ================= 连线逻辑 =================
    % 中心 -> 四大分支标题
    \draw[arrowStyle, red] (center) -- (value);
    \draw[arrowStyle, blue] (center) -- (batch);
    \draw[arrowStyle, green!60!black] (center) -- (sgd);
    \draw[arrowStyle, purple] (center) -- (mini);

    % 分支标题 -> 子节点 (上方)
    \draw[arrowStyle, red!60] (value) -- (iterative.south);
    \draw[arrowStyle, red!60] (value) -- (info.south);
    \draw[arrowStyle, red!60] (value) -- (complexity.south);
    \draw[arrowStyle, red!60] (value) -- (backprop.south);

    % 分支标题 -> 子节点 (右侧)
    \draw[arrowStyle, blue!60] (batch) -- (fullbatch.west);
    \draw[arrowStyle, blue!60] (batch) -- (lr.west);
    \draw[arrowStyle, blue!60] (batch) -- (steepest.west);
    \draw[arrowStyle, blue!60] (batch) -- (inefficient.west);

    % 分支标题 -> 子节点 (下方)
    \draw[arrowStyle, green!60!black] (sgd) -- (single.north);
    \draw[arrowStyle, green!60!black] (sgd) -- (epoch.north);
    \draw[arrowStyle, green!60!black] (sgd) -- (online.north);
    \draw[arrowStyle, green!60!black] (sgd) -- (redundancy.north);
    \draw[arrowStyle, green!60!black] (sgd) -- (escape.north);

    % 分支标题 -> 子节点 (左侧)
    \draw[arrowStyle, purple!60] (mini) -- (minibatch.east);
    \draw[arrowStyle, purple!60] (mini) -- (noise.east);
    \draw[arrowStyle, purple!60] (mini) -- (shuffle.east);
    \draw[arrowStyle, purple!60] (mini) -- (symmetry.east);
    \draw[arrowStyle, purple!60] (mini) -- (he.east);

\end{tikzpicture}

\end{document}